% Industry resume, Grover Aruquipa
% Compile with pdfLaTeX (also works on Overleaf). One column on purpose:
% applicant tracking systems read single-column text far more reliably than sidebars.
\documentclass[10pt,letterpaper]{article}
\usepackage[T1]{fontenc}
\usepackage[utf8]{inputenc}
\usepackage[scaled=0.95]{helvet}
\renewcommand{\familydefault}{\sfdefault}
\usepackage[top=1.0cm,bottom=0.8cm,left=1.5cm,right=1.5cm]{geometry}
\usepackage{enumitem}
\usepackage{titlesec}
\usepackage{xcolor}
\usepackage{microtype}
\usepackage[hidelinks]{hyperref}

\definecolor{ink}{HTML}{111111}
\definecolor{accent}{HTML}{0D2B61}
\definecolor{muted}{HTML}{555555}
\color{ink}
\pagestyle{empty}
\setlength{\parindent}{0pt}
\setlength{\parskip}{0pt}

\titleformat{\section}{\normalsize\bfseries\color{accent}}{}{0em}{\MakeUppercase}[\vspace{-6pt}{\color{accent}\rule{\linewidth}{0.6pt}}]
\titlespacing*{\section}{0pt}{5pt}{2pt}
\setlist[itemize]{leftmargin=1.1em,itemsep=0.5pt,topsep=1pt,parsep=0pt,label=\textbullet}

% \role{Title}{Organization}{Location}{Dates}
\newcommand{\role}[4]{\textbf{#1}, #2 \hfill {\color{muted}#3 \textbar\ #4}\par}

\begin{document}

%================ HEADER ================
{\LARGE\bfseries Grover Aruquipa}\par\vspace{2pt}
{\large Mechanical Engineer, Robotic Hand Design \textbar\ High-DoF Mechanisms, Remote Actuation, Tight Packaging}\par\vspace{3pt}
{\small Qu\'ebec, QC, Canada \ \textbullet\ \
\href{mailto:grover.aruquipa.9@gmail.com}{grover.aruquipa.9@gmail.com} \ \textbullet\ \
\href{https://www.linkedin.com/in/grover-aruquipa-996917126}{linkedin.com/in/grover-aruquipa-996917126} \ \textbullet\ \
\href{https://groveraruquipa.github.io}{groveraruquipa.github.io}}\par\vspace{1pt}
{\small\bfseries Available from January 2027 \ \textbullet\ \ Open to relocation}\par


%================ SUMMARY ================
\section{Summary}
Mechanical engineer finishing a PhD in robot mechanisms and grasping (Universit\'e Laval, defense early 2027), and
author of a published review of parallel mechanisms for grasping and robotic hands. I design high-DoF mechanisms whose
grippers are driven remotely, from motors at the base through linkages, concentric shafts and belts, and I take each one
from architecture and load paths through SolidWorks, FEA, prototyping and test. Three robots built, a fourth with an
integrated Sarrus-linkage gripper in progress.

%================ SKILLS ================
\section{Technical Skills}
\begin{itemize}
\item \textbf{Hand and mechanism design:} high-DoF parallel and hybrid mechanisms, integrated grippers, joint and load-path
architecture; remote actuation through linkages, concentric hollow shafts and belts; underactuated and parallel robot
hands (published review); compact packaging of actuation, encoders and wiring.
\item \textbf{CAD and analysis:} SolidWorks (expert), Creo; FEA in ANSYS for stress and stiffness; statics, kinematics,
grasp and force analysis, Jacobians and singularities; MATLAB, Simscape, MuJoCo, NVIDIA Isaac Sim, Pinocchio.
\item \textbf{Prototyping and production:} rapid prototyping (SLS, FDM, machining), DFM and DFA, fits and tolerances for
bearings and shafts, component and supplier selection; test setups for accuracy, repeatability and durability.
\item \textbf{Cross-functional:} real-time control at 500\,Hz (C++, Simulink), motor and encoder integration, ROS~2;
presented designs for review; technical lead of two graduate students.
\end{itemize}

%================ EXPERIENCE ================
\section{Experience}
\role{Doctoral Researcher}{Robotics Laboratory, Universit\'e Laval}{Qu\'ebec, Canada}{2023 -- present}
\begin{itemize}
\item \textbf{Designed, built and tested three high-DoF robots} (4, 7 and 8 actuators, all at the base) whose grippers are
actuated remotely through the mechanism; owned architecture, CAD, FEA, prototyping, integration and test.
\item \textbf{4-DoF coaxial wrist-gripper:} four actuation channels packaged in concentric hollow shafts, 180\,mm
footprint, 0.25\,kg moving mass; unlimited roll verified over ten revolutions with sub-degree tracking while grasping.
\item \textbf{Delta robot with integrated Sarrus-linkage gripper (in progress):} grasp force estimated from motor currents,
no sensor in the fingers, to pick fragile objects quickly.
\item \textbf{(6+2)-DoF robot:} 8 coaxial actuators, three orthogonal unlimited rotations and a reconfigurable grasping
platform; technical lead of two graduate students from concept to prototype.
\end{itemize}

\vspace{3pt}
\role{Visiting Researcher}{Institute of Robotics, Johannes Kepler University Linz}{Linz, Austria}{Apr. -- May 2026}
\begin{itemize}
\item Developed singularity and grasp-quality indices for manipulators with integrated grasping, hosted by Univ.-Prof.~A.~M\"uller (Mitacs Globalink Research Award).
\end{itemize}

\vspace{3pt}
\role{Research Intern}{FEMTO-ST (CNRS)}{Besan\c{c}on, France}{2022 -- 2023}
\begin{itemize}
\item Designed and modeled a redundant dexterous parallel robot with two unlimited rotations for micro-manipulation, and wrote its real-time C++ control software; clean-room fabrication projects.
\end{itemize}

\vspace{3pt}
\textbf{Industry roles} \hfill {\color{muted}2020 -- 2021}\par
{\small Computer Vision Engineer, Sencon Srl (industrial vision, scaffolding detection) \ \textbullet\ \ Data Science
Intern, Hitachi ABB Power Grids, Zurich (machine learning for high-voltage circuit breakers) \ \textbullet\ \ Data
Scientist, KTH Royal Institute of Technology.}

%================ EDUCATION ================
\section{Education}
\textbf{PhD, Mechanical Engineering}, Universit\'e Laval \hfill {\color{muted}Qu\'ebec, Canada \textbar\ 2023 -- 2027 (expected)}\par
{\small Kinematically redundant parallel robots with unlimited rotations and integrated grasping. Advisors: C.~Gosselin, P.~Lambert.}\par\vspace{2pt}
\textbf{MSc, Smart Integrated Systems}, Universit\'e Bourgogne--Franche-Comt\'e \hfill {\color{muted}Besan\c{c}on, France \textbar\ 2021 -- 2023}\par\vspace{2pt}
\textbf{BEng, Mechatronics Engineering}, Bolivian Catholic University \hfill {\color{muted}La Paz, Bolivia \textbar\ 2014 -- 2019}\par
{\small Thesis: reconfigurable Delta robot for product picking.}

%================ AWARDS ================
\section{Awards and Languages}
\begin{itemize}
\item 1st Place, Graduate Division, ASME Student Mechanisms and Robotics Design Competition (2025), team lead.
\item Finalist, MassRobotics competition, Boston, with one of my robot prototypes, team lead.
\item Mitacs Globalink Research Award (2026); 1st Place, Student Paper Competition, Argentine Robotics Conference (2019).
\item Languages: Spanish (native), English (full professional), French (intermediate).
\end{itemize}

\newpage
%================ PUBLICATIONS
\section{Selected Publications}
{\small First author on all papers. Full list at groveraruquipa.github.io.}\par\vspace{2pt}
\begin{itemize}
\item \textbf{G.~Aruquipa}, P.~Lambert, R.~Dahmouche, C.~Gosselin, ``Parallel robots for grasping tasks: a review,'' \emph{ASME Journal of Mechanisms and Robotics}, vol.~18, no.~12, art.~120801, 2026. doi:10.1115/1.4072612
\item \textbf{G.~Aruquipa}, P.~Lambert, C.~Gosselin, ``A compact coaxial parallel wrist-gripper with unlimited tool-axis rotation,'' under review, \emph{Robotics and Autonomous Systems}.
\item \textbf{G.~Aruquipa}, P.~Lambert, C.~Gosselin, ``Analysis and design of a novel (6+1)-DoF parallel robot with unlimited rotation and grasping capabilities,'' \emph{Mechanism and Machine Theory}, vol.~224, art.~106438, 2026.
\item \textbf{G.~Aruquipa}, P.~Lambert, C.~Gosselin, ``A kinematically redundant (6+2)-DoF parallel robot with three orthogonal unlimited rotations and coaxial actuation,'' submitted to \emph{IEEE ICRA 2027}.
\item \textbf{G.~Aruquipa}, V.~Muralidharan, P.~Lambert, C.~Gosselin, ``Analysis and design of a novel compact four-degree-of-freedom parallel robot with unlimited rotation,'' \emph{ASME Journal of Mechanisms and Robotics}, vol.~18, no.~6, art.~061001, 2026.
\item \textbf{G.~Aruquipa}, P.~Lambert, C.~Gosselin, ``Kinematic analysis and design of a novel 9-DoF parallel robot with grasping capabilities,'' \emph{CCToMM Symposium on Mechanisms, Machines and Mechatronics}, 2025, pp.~117--127.
\end{itemize}

%================ REFERENCES ================
\section{References}
\begin{itemize}
\item \textbf{Prof.~Cl\'ement Gosselin}, Full Professor, Department of Mechanical Engineering, Universit\'e Laval, Qu\'ebec,
Canada. PhD supervisor. \href{mailto:clement.gosselin@gmc.ulaval.ca}{clement.gosselin@gmc.ulaval.ca}
\item \textbf{Prof.~Patrice Lambert}, Assistant Professor, Department of Mechanical Engineering, Universit\'e Laval,
Qu\'ebec, Canada. PhD co-supervisor. \href{mailto:patrice.lambert@gmc.ulaval.ca}{patrice.lambert@gmc.ulaval.ca}
\item \textbf{Univ.-Prof.~Dr.-Ing.~Andreas M\"uller}, Head of the Institute of Robotics, Johannes Kepler University Linz,
Austria. Host during my research stay (2026). \href{mailto:a.mueller@jku.at}{a.mueller@jku.at}
\end{itemize}

\end{document}
